% Source: https://artofproblemsolving.com/community/c3h1263872_rhombus_properties
% Stable identity: aops-topic-1263872-q01
The diagonals of a rhombus bisect each other at right angles, cutting it into four congruent right triangles with legs $\dfrac{d_1}{2}$ and $\dfrac{d_2}{2}$.

\begin{asy}
usepackage("amsmath");
size(4cm);
defaultpen(fontsize(10pt));
pair A=(-2,0),B=(0,1.2),C=(2,0),D=(0,-1.2);
draw(A--B--C--D--cycle,linewidth(1)); draw(A--C,dashed); draw(B--D,dashed);
draw((0,0.15)--(0.15,0.15)--(0.15,0));
label("$\dfrac{d_1}{2}$",(-1,0),S); label("$\dfrac{d_2}{2}$",(0,0.6),W); label("$s$",(-1,0.6),NW);
label("\mbox{}",(min(currentpicture,user=true)+max(currentpicture,user=true))/2,black); // keeps the text after the figure black when LaTeX typesets the labels inline
\end{asy}

Each side is a hypotenuse: $s=\sqrt{\dfrac{d_1^2}{4}+\dfrac{d_2^2}{4}}=\dfrac{1}{2}\sqrt{d_1^2+d_2^2}$. So the perimeter is
\[4s=\boxed{2\sqrt{d_1^2+d_2^2}}.\]
